\documentclass[aspectratio=169,11pt]{beamer}
% =============================================================================
% Presentation 2 — Semester 2, end of Week 12 (11 April 2027)
%
% Follows Report 3 (Final, Sem 2 end of Week 8): present the completed
% project, its evaluation and your reflection, as in report3.tex. Replace
% the red \guide{...} prompts with your own content, and adjust the outline
% to the format and duration set by your Academic Advisor/Supervisor.
% =============================================================================
\newcommand{\presentationname}{Presentation 2}
% Shared preamble for the slides (ppt1.tex, ppt2.tex).
\usepackage[T1]{fontenc}
\usepackage[utf8]{inputenc}
\usepackage{booktabs} % For better-looking horizontal separators
\usepackage{graphicx}
\graphicspath{ {images/} }
\usepackage[
    backend=biber,
    style=apa,
    sorting=nyt
]{biblatex}
\addbibresource{cite.bib}
\usepackage{listings}
\lstset{
  basicstyle=\footnotesize\ttfamily,
  breaklines=true,
  frame=single
}

% Look: plain theme in NUS blue/orange, slide numbers, no navigation symbols
\definecolor{nusblue}{HTML}{003D7C}
\definecolor{nusorange}{HTML}{EF7C00}
\usecolortheme[named=nusblue]{structure}
\setbeamercolor{alerted text}{fg=nusorange}
\setbeamerfont{frametitle}{series=\bfseries}
\setbeamertemplate{navigation symbols}{}
\setbeamertemplate{footline}[frame number]
\setbeamertemplate{itemize items}[circle]
\setbeamertemplate{bibliography item}{}
\setbeamertemplate{frametitle continuation}{}

% Project configuration (course, title, members, advisor, repository)
% ---------------- Project configuration ----------------
% Shared by the report (main.tex) and the slides (ppt1.tex, ppt2.tex).
% Course: TCX3901 (BIT) or TIC3901 (BTech Computing)
\newcommand{\coursecode}{TCX3901}
% \newcommand{\coursecode}{TIC3901}
\newcommand{\coursename}{Industrial Practice}
\newcommand{\programme}{Bachelor of Information Technology Programme}
% \newcommand{\programme}{Bachelor of Technology Programme}

% Project type: Work-Related or Non-Work-Related (see README for the caveats)
\newcommand{\projectkind}{Non-Work-Related Project}
% \newcommand{\projectkind}{Work-Related Project}

% Individual (single-student) project: uncomment \individualprojecttrue. The
% contributions page then becomes a declaration of individual work.
\newif\ifindividualproject
% \individualprojecttrue

% List all group members, up to 5 (one per line with \\; delete unused lines),
% each with their own matriculation number. Every student must still submit
% the (same) file on ACES for tracking purposes.
\newcommand{\authorname}{%
    Member One (A0000001V)\\
    Member Two (A0000002W)\\
    Member Three (A0000003X)\\
    Member Four (A0000004Y)\\
    Member Five (A0000005Z)}
\newcommand{\advisorname}{Dr. Han Liang Wee, Eric}
\newcommand{\projecttitle}{Project Title}
\newcommand{\academicyear}{AY2026/27}
% Link to the group's source code repository. Create it early and add your
% supervisor (GitHub: eric-vader) as a collaborator so that it can be accessed.
\newcommand{\repourl}{https://github.com/your-group/your-repo}

% (No need to edit below.)
% Template guidance: \guide{...} marks prompts to the writer, in the report
% and the slides. They are typeset in red so that they are not mistaken for
% your own text. Replace all of them; none should remain in a submission.
\newcommand{\guide}[1]{{\itshape\color{red}#1}}

% The link is flagged in red until the placeholder above
% is replaced.
\newcommand{\placeholderrepourl}{https://github.com/your-group/your-repo}
\newcommand{\repolink}{\url{\repourl}%
    \ifx\repourl\placeholderrepourl
        \ \textcolor{red}{(replace and add supervisor to repo)}%
    \fi}



% \presentationname is defined in ppt1.tex / ppt2.tex
\title{\projecttitle}
\subtitle{\coursecode\ \coursename\ --- \presentationname}
\author[\coursecode]{\authorname}
\institute{Supervisor: \advisorname\\[1ex]
    \programme\\
    National University of Singapore}
\date{\academicyear\\[1ex]
    {\footnotesize Repository: \repolink}}


\begin{document}

\begin{frame}[plain]
    \titlepage
\end{frame}

\begin{frame}{Outline}
    \tableofcontents
\end{frame}

\section{Recap}

\begin{frame}{Problem and Objectives}
    \begin{itemize}
        \item \guide{Brief recap of the problem and motivation. Example citation
              \autocite{lamport_time_1978}.}
        \item \guide{The objectives set out in the proposal, and any changes to the
              objectives, scope or approach since Presentation 1.}
    \end{itemize}
\end{frame}

\section{Implementation}

\begin{frame}{Implementation}
    \begin{itemize}
        \item \guide{What was built and delivered.}
        \item \guide{How it fulfils the objectives.}
    \end{itemize}
\end{frame}

\begin{frame}{Design Overview}
    \begin{itemize}
        \item \guide{The final design or architecture (a diagram from
              \texttt{images/} works best).}
        \item \guide{Key decisions and their trade-offs.}
    \end{itemize}
\end{frame}

\begin{frame}{Demo}
    \begin{itemize}
        \item \guide{Live demo or screenshots of the delivered project.}
    \end{itemize}
\end{frame}

\section{Evaluation}

\begin{frame}{Evaluation}
    \begin{itemize}
        \item \guide{How the outcome was evaluated: methodology, setup and metrics.}
    \end{itemize}
\end{frame}

\begin{frame}{Results}
    \begin{itemize}
        \item \guide{Results and their interpretation, related back to the
              objectives and user stories.}
    \end{itemize}
\end{frame}

\begin{frame}{Challenges}
    \begin{itemize}
        \item \guide{Difficulties encountered and how they were mitigated.}
    \end{itemize}
\end{frame}

\section{Conclusion}

\begin{frame}{Summary}
    \begin{itemize}
        \item \guide{What was achieved, measured against the objectives.}
    \end{itemize}
\end{frame}

% This reflection is the core of the course objective: how concepts and
% theories gained in the classroom translate into (industrial) practice.
\begin{frame}{Reflection}
    \begin{itemize}
        \item \guide{How concepts and theories from your degree programme were
              applied in the project, and how they affected its outcome.}
    \end{itemize}
\end{frame}

\begin{frame}{Future Work}
    \begin{itemize}
        \item \guide{Extensions or improvements that could follow from this project.}
    \end{itemize}
\end{frame}

% Group projects only (hidden when \individualprojecttrue is set in config.tex).
% Delete the rows for unused members, as in contributions.tex.
\ifindividualproject\else
\begin{frame}{Individual Contributions}
    \centering
    \begin{tabular}{llp{0.4\textwidth}}
        \toprule
        Member & Role & Contributions \\
        \midrule
        Member One & Team Lead & \\
        Member Two & Backend Developer & \\
        Member Three & Frontend Developer & \\
        Member Four & QA \& Testing Lead & \\
        Member Five & DevOps \& Documentation & \\
        \bottomrule
    \end{tabular}

    \vspace{1em}
    \guide{\small Not needed for an individual project: set
    \texttt{\textbackslash individualprojecttrue} in \texttt{config.tex} to
    remove this slide.}
\end{frame}
\fi

\begin{frame}[allowframebreaks]{References}
    \printbibliography[heading=none]
\end{frame}

\begin{frame}[plain]
    \centering
    {\Large\bfseries\usebeamercolor[fg]{structure} Thank you\par}
    \vspace{1em}
    Questions?\par
    \vspace{1em}
    {\small\url{\repourl}}
\end{frame}

\end{document}
